% Options for packages loaded elsewhere
% Options for packages loaded elsewhere
\PassOptionsToPackage{unicode}{hyperref}
\PassOptionsToPackage{hyphens}{url}
\PassOptionsToPackage{dvipsnames,svgnames,x11names}{xcolor}
%
\documentclass[
  letterpaper,
  DIV=11,
  numbers=noendperiod]{scrartcl}
\usepackage{xcolor}
\usepackage{amsmath,amssymb}
\setcounter{secnumdepth}{-\maxdimen} % remove section numbering
\usepackage{iftex}
\ifPDFTeX
  \usepackage[T1]{fontenc}
  \usepackage[utf8]{inputenc}
  \usepackage{textcomp} % provide euro and other symbols
\else % if luatex or xetex
  \usepackage{unicode-math} % this also loads fontspec
  \defaultfontfeatures{Scale=MatchLowercase}
  \defaultfontfeatures[\rmfamily]{Ligatures=TeX,Scale=1}
\fi
\usepackage{lmodern}
\ifPDFTeX\else
  % xetex/luatex font selection
\fi
% Use upquote if available, for straight quotes in verbatim environments
\IfFileExists{upquote.sty}{\usepackage{upquote}}{}
\IfFileExists{microtype.sty}{% use microtype if available
  \usepackage[]{microtype}
  \UseMicrotypeSet[protrusion]{basicmath} % disable protrusion for tt fonts
}{}
\makeatletter
\@ifundefined{KOMAClassName}{% if non-KOMA class
  \IfFileExists{parskip.sty}{%
    \usepackage{parskip}
  }{% else
    \setlength{\parindent}{0pt}
    \setlength{\parskip}{6pt plus 2pt minus 1pt}}
}{% if KOMA class
  \KOMAoptions{parskip=half}}
\makeatother
% Make \paragraph and \subparagraph free-standing
\makeatletter
\ifx\paragraph\undefined\else
  \let\oldparagraph\paragraph
  \renewcommand{\paragraph}{
    \@ifstar
      \xxxParagraphStar
      \xxxParagraphNoStar
  }
  \newcommand{\xxxParagraphStar}[1]{\oldparagraph*{#1}\mbox{}}
  \newcommand{\xxxParagraphNoStar}[1]{\oldparagraph{#1}\mbox{}}
\fi
\ifx\subparagraph\undefined\else
  \let\oldsubparagraph\subparagraph
  \renewcommand{\subparagraph}{
    \@ifstar
      \xxxSubParagraphStar
      \xxxSubParagraphNoStar
  }
  \newcommand{\xxxSubParagraphStar}[1]{\oldsubparagraph*{#1}\mbox{}}
  \newcommand{\xxxSubParagraphNoStar}[1]{\oldsubparagraph{#1}\mbox{}}
\fi
\makeatother

\usepackage{color}
\usepackage{fancyvrb}
\newcommand{\VerbBar}{|}
\newcommand{\VERB}{\Verb[commandchars=\\\{\}]}
\DefineVerbatimEnvironment{Highlighting}{Verbatim}{commandchars=\\\{\}}
% Add ',fontsize=\small' for more characters per line
\usepackage{framed}
\definecolor{shadecolor}{RGB}{241,243,245}
\newenvironment{Shaded}{\begin{snugshade}}{\end{snugshade}}
\newcommand{\AlertTok}[1]{\textcolor[rgb]{0.68,0.00,0.00}{#1}}
\newcommand{\AnnotationTok}[1]{\textcolor[rgb]{0.37,0.37,0.37}{#1}}
\newcommand{\AttributeTok}[1]{\textcolor[rgb]{0.40,0.45,0.13}{#1}}
\newcommand{\BaseNTok}[1]{\textcolor[rgb]{0.68,0.00,0.00}{#1}}
\newcommand{\BuiltInTok}[1]{\textcolor[rgb]{0.00,0.23,0.31}{#1}}
\newcommand{\CharTok}[1]{\textcolor[rgb]{0.13,0.47,0.30}{#1}}
\newcommand{\CommentTok}[1]{\textcolor[rgb]{0.37,0.37,0.37}{#1}}
\newcommand{\CommentVarTok}[1]{\textcolor[rgb]{0.37,0.37,0.37}{\textit{#1}}}
\newcommand{\ConstantTok}[1]{\textcolor[rgb]{0.56,0.35,0.01}{#1}}
\newcommand{\ControlFlowTok}[1]{\textcolor[rgb]{0.00,0.23,0.31}{\textbf{#1}}}
\newcommand{\DataTypeTok}[1]{\textcolor[rgb]{0.68,0.00,0.00}{#1}}
\newcommand{\DecValTok}[1]{\textcolor[rgb]{0.68,0.00,0.00}{#1}}
\newcommand{\DocumentationTok}[1]{\textcolor[rgb]{0.37,0.37,0.37}{\textit{#1}}}
\newcommand{\ErrorTok}[1]{\textcolor[rgb]{0.68,0.00,0.00}{#1}}
\newcommand{\ExtensionTok}[1]{\textcolor[rgb]{0.00,0.23,0.31}{#1}}
\newcommand{\FloatTok}[1]{\textcolor[rgb]{0.68,0.00,0.00}{#1}}
\newcommand{\FunctionTok}[1]{\textcolor[rgb]{0.28,0.35,0.67}{#1}}
\newcommand{\ImportTok}[1]{\textcolor[rgb]{0.00,0.46,0.62}{#1}}
\newcommand{\InformationTok}[1]{\textcolor[rgb]{0.37,0.37,0.37}{#1}}
\newcommand{\KeywordTok}[1]{\textcolor[rgb]{0.00,0.23,0.31}{\textbf{#1}}}
\newcommand{\NormalTok}[1]{\textcolor[rgb]{0.00,0.23,0.31}{#1}}
\newcommand{\OperatorTok}[1]{\textcolor[rgb]{0.37,0.37,0.37}{#1}}
\newcommand{\OtherTok}[1]{\textcolor[rgb]{0.00,0.23,0.31}{#1}}
\newcommand{\PreprocessorTok}[1]{\textcolor[rgb]{0.68,0.00,0.00}{#1}}
\newcommand{\RegionMarkerTok}[1]{\textcolor[rgb]{0.00,0.23,0.31}{#1}}
\newcommand{\SpecialCharTok}[1]{\textcolor[rgb]{0.37,0.37,0.37}{#1}}
\newcommand{\SpecialStringTok}[1]{\textcolor[rgb]{0.13,0.47,0.30}{#1}}
\newcommand{\StringTok}[1]{\textcolor[rgb]{0.13,0.47,0.30}{#1}}
\newcommand{\VariableTok}[1]{\textcolor[rgb]{0.07,0.07,0.07}{#1}}
\newcommand{\VerbatimStringTok}[1]{\textcolor[rgb]{0.13,0.47,0.30}{#1}}
\newcommand{\WarningTok}[1]{\textcolor[rgb]{0.37,0.37,0.37}{\textit{#1}}}

\usepackage{longtable,booktabs,array}
\newcounter{none} % for unnumbered tables
\usepackage{calc} % for calculating minipage widths
% Correct order of tables after \paragraph or \subparagraph
\usepackage{etoolbox}
\makeatletter
\patchcmd\longtable{\par}{\if@noskipsec\mbox{}\fi\par}{}{}
\makeatother
% Allow footnotes in longtable head/foot
\IfFileExists{footnotehyper.sty}{\usepackage{footnotehyper}}{\usepackage{footnote}}
\makesavenoteenv{longtable}
\usepackage{graphicx}
\makeatletter
\newsavebox\pandoc@box
\newcommand*\pandocbounded[1]{% scales image to fit in text height/width
  \sbox\pandoc@box{#1}%
  \Gscale@div\@tempa{\textheight}{\dimexpr\ht\pandoc@box+\dp\pandoc@box\relax}%
  \Gscale@div\@tempb{\linewidth}{\wd\pandoc@box}%
  \ifdim\@tempb\p@<\@tempa\p@\let\@tempa\@tempb\fi% select the smaller of both
  \ifdim\@tempa\p@<\p@\scalebox{\@tempa}{\usebox\pandoc@box}%
  \else\usebox{\pandoc@box}%
  \fi%
}
% Set default figure placement to htbp
\def\fps@figure{htbp}
\makeatother





\setlength{\emergencystretch}{3em} % prevent overfull lines

\providecommand{\tightlist}{%
  \setlength{\itemsep}{0pt}\setlength{\parskip}{0pt}}



 


\KOMAoption{captions}{tableheading}
\makeatletter
\@ifpackageloaded{caption}{}{\usepackage{caption}}
\AtBeginDocument{%
\ifdefined\contentsname
  \renewcommand*\contentsname{Table of contents}
\else
  \newcommand\contentsname{Table of contents}
\fi
\ifdefined\listfigurename
  \renewcommand*\listfigurename{List of Figures}
\else
  \newcommand\listfigurename{List of Figures}
\fi
\ifdefined\listtablename
  \renewcommand*\listtablename{List of Tables}
\else
  \newcommand\listtablename{List of Tables}
\fi
\ifdefined\figurename
  \renewcommand*\figurename{Figure}
\else
  \newcommand\figurename{Figure}
\fi
\ifdefined\tablename
  \renewcommand*\tablename{Table}
\else
  \newcommand\tablename{Table}
\fi
}
\@ifpackageloaded{float}{}{\usepackage{float}}
\floatstyle{ruled}
\@ifundefined{c@chapter}{\newfloat{codelisting}{h}{lop}}{\newfloat{codelisting}{h}{lop}[chapter]}
\floatname{codelisting}{Listing}
\newcommand*\listoflistings{\listof{codelisting}{List of Listings}}
\makeatother
\makeatletter
\makeatother
\makeatletter
\@ifpackageloaded{caption}{}{\usepackage{caption}}
\@ifpackageloaded{subcaption}{}{\usepackage{subcaption}}
\makeatother
\usepackage{bookmark}
\IfFileExists{xurl.sty}{\usepackage{xurl}}{} % add URL line breaks if available
\urlstyle{same}
% fallback for those not using the hyperref driver hyperxmp:
\makeatletter
\@ifundefined{xmpquote}{\newcommand{\xmpquote}[1]{#1}}{}
\makeatother
\hypersetup{
  colorlinks=true,
  linkcolor={blue},
  filecolor={Maroon},
  citecolor={Blue},
  urlcolor={Blue},
  pdfcreator={LaTeX via pandoc}}


\author{}
\date{}
\begin{document}


\section{\texorpdfstring{\texttt{spurious\_stationary\_sim\_refactored.R}:
What Changed and
Why}{spurious\_stationary\_sim\_refactored.R: What Changed and Why}}\label{spurious_stationary_sim_refactored.r-what-changed-and-why}

\subsection{Scope}\label{scope}

This note explains the refactor from
\texttt{scripts/spurious\_stationary\_sim.R} to
\texttt{scripts/spurious\_stationary\_sim\_refactored.R}. It covers code
structure only. No simulation logic, formula, or random-number call
changed, and none of the theoretical content in
\texttt{spurious\_regression\_stationary.qmd} is affected.

\subsection{Guarantee: behaviour is
unchanged}\label{guarantee-behaviour-is-unchanged}

Every call that consumes randomness (\texttt{rnorm}, \texttt{draw\_ar1},
\texttt{draw\_ma2}, and the \texttt{ar\_pair}/\texttt{ma2\_pair} draw
closures built from them) appears in the refactored script in the same
order, with the same arguments, as in the original. Nothing was
reordered, added, or removed from the random-number stream.
Consequently, run with the same seeds, the refactored script reproduces
the original's CSV and figure output byte-for-byte. The change is
confined to eliminating duplicated bookkeeping code that computes no
randomness itself.

\subsection{The four changes}\label{the-four-changes}

\subsubsection{\texorpdfstring{1. \texttt{ar\_pair\_dgp()} replaces six
repeated registry
entries}{1. ar\_pair\_dgp() replaces six repeated registry entries}}\label{ar_pair_dgp-replaces-six-repeated-registry-entries}

The original \texttt{dgp\_registry} list built each of the six
AR(1)/AR(1) pairs ---\texttt{(0.5,0.5)}, \texttt{(0.8,0.8)},
\texttt{(0.9,0.9)}, \texttt{(0.95,0.95)}, \texttt{(0.9,0.5)},
\texttt{(0.8,-0.8)}---by hand, each one a four-line
\texttt{list(dgp\ =\ ...,\ type\ =\ ...,\ lambda\ =\ lambda\_ar\_pair(...),\ draw\ =\ ar\_pair(...))}
block that repeated \texttt{phi\_x} and \texttt{phi\_y} twice. The
refactored script factors this into one constructor:

\begin{Shaded}
\begin{Highlighting}[]
\NormalTok{ar\_pair\_dgp }\OtherTok{\textless{}{-}} \ControlFlowTok{function}\NormalTok{(name, phi\_x, phi\_y) \{}
  \FunctionTok{list}\NormalTok{(}\AttributeTok{dgp =}\NormalTok{ name, }\AttributeTok{type =} \StringTok{"stationary"}\NormalTok{,}
       \AttributeTok{lambda =} \FunctionTok{lambda\_ar\_pair}\NormalTok{(phi\_x, phi\_y),}
       \AttributeTok{draw =} \FunctionTok{ar\_pair}\NormalTok{(phi\_x, phi\_y))}
\NormalTok{\}}
\end{Highlighting}
\end{Shaded}

and each registry entry becomes a single line,
e.g.~\texttt{ar\_pair\_dgp("(0.9,0.5)",\ 0.9,\ 0.5)}.
\texttt{lambda\_ar\_pair()} and \texttt{ar\_pair()} are still called
with the same arguments in the same order as before;
\texttt{ar\_pair\_dgp()} only removes the need to write that pair of
calls out longhand six times.

\subsubsection{\texorpdfstring{2. \texttt{asymptotic\_rejection()}
replaces a formula written out
twice}{2. asymptotic\_rejection() replaces a formula written out twice}}\label{asymptotic_rejection-replaces-a-formula-written-out-twice}

The asymptotic rejection probability implied by a given long-run
variance ratio,

{[} 2,\Phi!\left(-,z\_\{\alpha/2\}/\sqrt{\lambda}\right), {]}

was written out independently in two places: once to build the
\texttt{asymptotic} column of \texttt{dgp\_spec}, and again (reformatted
across two lines) for \texttt{common\_volatility\_asymptotic}. The
refactored script defines it once:

\begin{Shaded}
\begin{Highlighting}[]
\NormalTok{asymptotic\_rejection }\OtherTok{\textless{}{-}} \ControlFlowTok{function}\NormalTok{(lambda, alpha) \{}
  \DecValTok{2} \SpecialCharTok{*} \FunctionTok{pnorm}\NormalTok{(}\SpecialCharTok{{-}}\FunctionTok{qnorm}\NormalTok{(}\DecValTok{1} \SpecialCharTok{{-}}\NormalTok{ alpha }\SpecialCharTok{/} \DecValTok{2}\NormalTok{) }\SpecialCharTok{/} \FunctionTok{sqrt}\NormalTok{(lambda))}
\NormalTok{\}}
\end{Highlighting}
\end{Shaded}

Both call sites now read
\texttt{asymptotic\_rejection(dgp\_spec\$lambda,\ alpha)} and
\texttt{asymptotic\_rejection(common\_volatility\_lambda,\ alpha)}
respectively. The \texttt{ifelse(is.na(lambda),\ NA\_real\_,\ ...)}
guard in the original's first call site is unnecessary --- \texttt{NA}
lambda propagates to \texttt{NA} through arithmetic and \texttt{pnorm()}
on its own --- so the refactored version relies on that instead of
guarding explicitly; this is noted in a comment at the definition site.

\subsubsection{\texorpdfstring{3. \texttt{with\_mcse()} replaces six
repeated Monte Carlo standard-error
lines}{3. with\_mcse() replaces six repeated Monte Carlo standard-error lines}}\label{with_mcse-replaces-six-repeated-monte-carlo-standard-error-lines}

Six separate result tables (\texttt{res}, \texttt{hac\_res},
\texttt{common\_volatility\_res}, \texttt{local\_to\_unity\_res},
\texttt{granger\_res}, \texttt{power\_res}) each had a line of the form
\texttt{\textless{}table\textgreater{}\$mcse\ \textless{}-\ mcse(\textless{}table\textgreater{}\$\textless{}column\textgreater{},\ \textless{}table\textgreater{}\$Nsim)}
attaching the Monte Carlo standard error of a reported proportion. Five
use the \texttt{reject} column; \texttt{power\_res} uses \texttt{power}.
The refactored script factors the pattern into:

\begin{Shaded}
\begin{Highlighting}[]
\NormalTok{with\_mcse }\OtherTok{\textless{}{-}} \ControlFlowTok{function}\NormalTok{(df, }\AttributeTok{value\_col =} \StringTok{"reject"}\NormalTok{) \{}
\NormalTok{  df}\SpecialCharTok{$}\NormalTok{mcse }\OtherTok{\textless{}{-}} \FunctionTok{mcse}\NormalTok{(df[[value\_col]], df}\SpecialCharTok{$}\NormalTok{Nsim)}
\NormalTok{  df}
\NormalTok{\}}
\end{Highlighting}
\end{Shaded}

and each of the six call sites becomes
\texttt{res\ \textless{}-\ with\_mcse(res)} (or, for the one table with
a differently named column,
\texttt{power\_res\ \textless{}-\ with\_mcse(power\_res,\ "power")}).
The underlying \texttt{mcse()} function is untouched.

\subsubsection{\texorpdfstring{4. Removed
\texttt{stationary\_spec}}{4. Removed stationary\_spec}}\label{removed-stationary_spec}

\texttt{stationary\_spec\ \textless{}-\ subset(dgp\_spec,\ type\ ==\ "stationary")}
was computed immediately after \texttt{dgp\_spec} but never read
anywhere else in the script. It was deleted.

\subsection{What did not change}\label{what-did-not-change}

\begin{itemize}
\tightlist
\item
  Every DGP definition, sample-size grid, seed, and replication count.
\item
  Every simulation function (\texttt{simulate\_dgp},
  \texttt{simulate\_common\_volatility},
  \texttt{simulate\_local\_to\_unity}, \texttt{simulate\_granger},
  \texttt{simulate\_power}) and its internal logic.
\item
  \texttt{mcse()}, \texttt{lambda\_ar\_pair()}, \texttt{ar\_pair()},
  \texttt{draw\_ar1()}, \texttt{draw\_ma2()}, and every other existing
  helper.
\item
  All file output: CSV paths, plot code, and figure output.
\end{itemize}

\subsection{Validation performed}\label{validation-performed}

Both scripts were run end to end in quick mode
(\texttt{SPURIOUS\_QUICK=1}, same seeds) to separate output directories,
and all nine CSV outputs (\texttt{spurious\_common\_volatility.csv},
\texttt{spurious\_comparison\_table.csv},
\texttt{spurious\_granger\_misspecification.csv},
\texttt{spurious\_hac\_comparison\_table.csv},
\texttt{spurious\_hac\_rejection\_frequency\_series.csv},
\texttt{spurious\_local\_power\_critical\_values.csv},
\texttt{spurious\_local\_power.csv},
\texttt{spurious\_local\_to\_unity.csv},
\texttt{spurious\_rejection\_frequency\_series.csv}) were confirmed
byte-identical between the two runs. This is consistent with the
behaviour-preserving guarantee above; it was not re-checked under the
full (non-quick) 2,000-replication run, but the random-number call
sequence is identical between the two scripts regardless of replication
count, so the same equivalence is expected to hold there too.

\subsection{Outcome}\label{outcome}

The refactor is purely a readability and maintainability change: three
call sites of duplicated logic (registry construction, the asymptotic
rejection formula, and Monte Carlo standard-error attachment) are now
each expressed once, and one dead variable was removed. The simulation
design, its outputs, and the theoretical document it supports are
unaffected.




\end{document}
